\chapter{Agent 的构成}
\chapterlead{命令行里的“AI 助手”不是聊天窗口，而是一个会调用工具的循环程序。本章拆开这个循环，说明它为什么能改文件、会犯什么错，以及权限为什么必须由外部约束。}

\section{起步：观察一次工具调用}

装一个最小的 Agent 之前，先理解它在做什么。一次典型交互包含五步。

\begin{booktable}{\textwidth}{J{0.667} J{1.333}}
\toprule
步骤 & 发生的事 \\
\midrule
一 & 用户提出目标，附带当前环境信息 \\
二 & 模型输出一段“要调用工具”的结构化请求 \\
三 & 本地程序执行该工具，把结果回填给模型 \\
四 & 模型据结果决定下一步，或给出结论 \\
五 & 全部对话与工具输出写入会话文件 \\
\bottomrule
\end{booktable}

关键在于第二步与第三步之间：模型本身不能读写文件，它只能输出“请执行这个命令”。真正执行的是本地程序。因此安全边界在程序一侧，不在模型一侧。

\begin{figure}[H]
\centering
\figfit{%
\begin{tikzpicture}[font=\small, node distance=8mm,
  box/.style={draw=ink!45, fill=ink!5, rounded corners=1pt, inner sep=4pt, align=center, minimum height=9mm},
  arr/.style={-{Latex[length=1.5mm]}, teal, thick}]
  \node[box] (user) {目标};
  \node[box, right=of user] (model) {模型\\[-1pt]{\scriptsize 只输出文本}};
  \node[box, right=of model] (tools) {工具\\[-1pt]{\scriptsize 读写、执行}};
  \node[box, below=of tools] (fs) {文件系统与\\进程};
  \draw[arr] (user) -- (model);
  \draw[arr] (model) -- node[above, font=\scriptsize] {调用请求} (tools);
  \draw[arr] (tools) -- (fs);
  \draw[arr, dashed] (tools.south west) -- node[below, font=\scriptsize] {结果回填} (model.south east);
\end{tikzpicture}%
}
\caption{工具调用循环。虚线是回填路径，它让模型看到执行结果}
\end{figure}

\section{工具有哪几类}

不同 Agent 的内置工具大同小异，通常有六类。

\begin{booktable}{\textwidth}{J{0.741} J{1.259}}
\toprule
类别 & 典型能力 \\
\midrule
读 & 查看文件、列目录、搜索内容 \\
写 & 新建、修改、删除文件 \\
执行 & 运行命令、跑测试、构建 \\
检索 & 按正则或语义在网络与本地查找 \\
任务 & 拆分待办、维护计划 \\
扩展 & 调用 MCP 服务、加载技能说明 \\
\bottomrule
\end{booktable}

执行类工具风险最高：一条命令可能删除文件、重装系统、向外发送数据。写类工具次之：改动可能覆盖未提交的工作。读类工具几乎无害，只是消耗上下文。

\section{上下文与压缩}

模型的输入窗口有限。一次会话里反复读文件、跑命令，很快就会填满。

\begin{shellsession}
$ wc -l ~/.omp/agent/sessions/*/*.jsonl | tail -1
$ du -sh ~/.omp/agent/sessions
\end{shellsession}

会话以 JSONL 保存，一行一条消息，包含工具调用与输出。长会话会触发压缩：把较早的对话摘要成一段，保留最近的原文。压缩是有损的，因此重要结论应让 Agent 写进文件，而不是留在对话里。

\section{为什么会犯错}

四类错误的成因不同，处理方式也不同。

\begin{enumerate}
  \item \textbf{信息不足}：没读到关键文件就下结论。表现为改错地方、漏掉配置。
  \item \textbf{过度自信}：把猜测写成事实。表现为引用不存在的函数、编造路径。
  \item \textbf{局部最优}：只改看得见的一处，不追调用方。表现为修好一个用例、弄坏另一处。
  \item \textbf{权限过大}：能做的事超出必要范围。表现为一条命令造成大范围改动。
\end{enumerate}

前两类靠提供上下文与要求验证；第三类靠测试与静态检查；第四类只能靠外部约束，也就是本书第三十四章的主题。

\section{Agent 与脚本的差别}

脚本按固定顺序执行，Agent 按目标自行决定步骤。这个差别决定了使用方式。

\begin{booktable}{\textwidth}{J{0.667} J{1.333}}
\toprule
维度 & 差别 \\
\midrule
可预测性 & 脚本确定，Agent 每次路径可能不同 \\
适用任务 & 步骤明确用脚本，需要判断用 Agent \\
验证方式 & 脚本看输出，Agent 看结果与差异 \\
成本 & 脚本几乎为零，Agent 按 token 计费 \\
\bottomrule
\end{booktable}

判断标准很简单：如果任务能用十行脚本写清楚，就写脚本；如果需要读若干文件、权衡几种做法、边做边验证，才交给 Agent。

\section{练习}

\begin{enumerate}
  \item 画出一次你使用过的 AI 交互的五个步骤，标出哪一步在本地执行。
  \item 列出你希望 Agent 拥有的三类工具，说明每一类可能造成的最大损失。
  \item 找到本机某个 CLI 工具的会话目录，统计一个会话文件的行数与大小。
  \item 举一个“应当写脚本而不是交给 Agent”的例子，说明理由。
  \item 举一个“适合交给 Agent”的例子，说明它需要哪些工具权限。
\end{enumerate}

\begin{checkpoint}
本章结束时，应能：
\begin{itemize}
  \item 描述工具调用循环的五个步骤，指出安全边界在哪；
  \item 按风险排序六类工具；
  \item 说明上下文压缩的代价，知道重要结论应落到文件；
  \item 区分四类错误并说出对应的处理方式；
  \item 判断一项任务该写脚本还是交给 Agent。
\end{itemize}
\end{checkpoint}
